% Purpose: Appendix H, the latent probe on the network: protocol, the two-part rule, the controls and the outcome
% (Figure H1). Converted from the Supplementary Materials, Section S7 (paper/si/sections/s7_probe.tex), and the former
% Methods section on the probe (sec:methods-probe, before R3), in session A4, with the network's probe record only.
% Input by main.tex after \appendix. Numbers: results/probe_descriptive.json (block s_v2; \numNetProbe... macros),
% configs/config.yaml (probe).
\section{Latent Probe}
\label{app:probe}

If learning formation energies from structure alone taught the network the electronic structure of the host, a linear map from its learned description of the vacancy site and its host to the DFT descriptors should improve as the network sees more training hosts, and should do better than the same map applied to an untrained network reading the same geometry. For every resample and training size, the readout input of the first-seed S network (the vacancy-site vector, the pooled host vector and the global state) is mapped by a ridge regression \citep{Pedregosa2018scikitlearn}, a linear least-squares fit with a penalty on the size of its coefficients, to each descriptor given to D \citep{Alain2018understanding}. The penalty is chosen by a \numCfgProbeAlphaCvFolds{}-fold cross-validation on the training sites grouped by host, and the probe is scored by the coefficient of determination \Rsq{} on the test sites, per descriptor and as the mean over each descriptor class.

The rule has two conditions for each class. The trend condition asks that the recovered fraction of descriptor variance rise with the number of training hosts: the Spearman correlation between the logarithm of the training size and the class-mean \Rsq{} over all \numProbeDescriptiveSVTwoHostRuleTrendNCells{} resample-size cells must have an interval above zero. The control condition asks that, at every training size from \numCfgProbePassRuleBeatsControlsMinBudget{} hosts on, the trained probe exceed two controls, each difference with an interval above zero: the same probe fitted to descriptor rows permuted among the training hosts (shuffled labels), and the probe on an untrained network of the same architecture and settings.

The trend condition holds for both classes (Figure~\ref{fig:app_probe}), with Spearman correlations of \numNetProbeHostTrendRho{} [\numNetProbeHostTrendRhoLo{}, \numNetProbeHostTrendRhoHi{}] for the host-electronic and \numNetProbeSiteTrendRho{} [\numNetProbeSiteTrendRhoLo{}, \numNetProbeSiteTrendRhoHi{}] for the site-electronic class; the class-mean \Rsq{} of the host-electronic descriptors rises from \numNetProbeHostRsqTwoHundred{} with \numCfgBudgetsHostsItemFour{} training hosts to \numNetProbeHostRsqSixHundredFiftyFour{} with \numCfgBudgetsMax{}. The trained probe exceeds the shuffled-label control at every size the rule names, by \numNetProbeHostMinusShuffledSixHundredFiftyFour{} [\numNetProbeHostMinusShuffledSixHundredFiftyFourLo{}, \numNetProbeHostMinusShuffledSixHundredFiftyFourHi{}] for the host class with \numCfgBudgetsMax{} training hosts. It does not exceed the untrained network for the host class at any of them: the difference is \numNetProbeHostMinusUntrainedTwoHundred{} [\numNetProbeHostMinusUntrainedTwoHundredLo{}, \numNetProbeHostMinusUntrainedTwoHundredHi{}] with \numCfgBudgetsHostsItemFour{} and \numNetProbeHostMinusUntrainedSixHundredFiftyFour{} [\numNetProbeHostMinusUntrainedSixHundredFiftyFourLo{}, \numNetProbeHostMinusUntrainedSixHundredFiftyFourHi{}] with \numCfgBudgetsMax{} training hosts. For the site class the trained probe lies below the untrained one with \numCfgBudgetsHostsItemFour{} training hosts (\numNetProbeSiteMinusUntrainedTwoHundred{} [\numNetProbeSiteMinusUntrainedTwoHundredLo{}, \numNetProbeSiteMinusUntrainedTwoHundredHi{}]) and above it only with \numCfgBudgetsMax{} (\numNetProbeSiteMinusUntrainedSixHundredFiftyFour{} [\numNetProbeSiteMinusUntrainedSixHundredFiftyFourLo{}, \numNetProbeSiteMinusUntrainedSixHundredFiftyFourHi{}]). Over all descriptors, the class-mean \Rsq{} with \numCfgBudgetsMax{} training hosts is \numNetProbeAllRsqSixHundredFiftyFour{} for the trained network and \numProbeDescriptiveSVTwoAllCurvesInitZeroSixHundredFiftyFourMean{} for the untrained one. The control condition therefore fails for both classes, and training on formation energies adds no linearly recoverable host-electronic information beyond what an untrained encoder of the same structure already exposes (Section~\ref{sec:disc_principal}). The probe measures linear recoverability, a lower bound on what the vectors encode, because the readout that maps them to the energy is nonlinear.

\begin{figure}[H]
\centering
\figph{figureH1_probe}
\caption{Linear recovery of the DFT descriptors from S. A ridge probe maps the readout input of S to each descriptor of D, scored by test \Rsq{} for the trained network, shuffled labels and an untrained encoder, as (\textbf{a}) host-electronic and (\textbf{b}) site-electronic class means against training hosts and (\textbf{c}) per descriptor at the largest training size. Bands, clipped at the axes, and bars are two-sided resample-bootstrap intervals, values in \texttt{results/probe\_descriptive.json}.\label{fig:app_probe}}
\end{figure}
